\section{Number of Agents}\label{sec:agents}
%% FUNCTION: Shows that with three goods per agent and remainder zero (m = 6, 9, 12, 15), EFX
%% becomes rarer as more agents take part, and that failures become deeper with agreement but not
%% without it. There is no special case with a proof. Bridge to Chapter 5.

The number of agents is the number of people who share. Two siblings who divide six items and
five siblings who divide fifteen items both pick three times each, but the five have many more
pairs among them who might envy each other. Because $n$, $m$, $r$ and $q$ are linked
(Section~\ref{sec:design}), we compare different numbers of agents only at the same number of
goods per agent: we keep $r=3$ and $q=0$ and vary $n$ from $2$ to $5$, which means $m=6$, $9$, $12$
and $15$ goods. Agreement and dispersion stay at the reference levels, Spearman $0.48$
($\rho=0.5$) and Gini $0.52$ ($\theta=1$).

\subsection{EFX Rate}\label{sec:agents-rate}

As Figure~\ref{fig:agents} shows, Round Robin achieves EFX less often when more agents share,
even though every agent receives the same number of goods. With two agents, the EFX rate is
$0.687$, with $6{,}256$ failures among $R=20{,}000$ instances. It falls to $0.459$ with three
agents, $0.293$ with four and $0.188$ with five, where $16{,}231$ instances fail (measured). Each
additional agent costs about a third of the remaining EFX rate.

\begin{figure}[ht]
  \centering
  \includegraphics{fig45_agents_rate.pdf}
  \caption{EFX rate of Round Robin against the number of agents $n$, with $m=3n$ goods ($r=3$,
  $q=0$), Spearman $0.48$ ($\rho=0.5$) and Gini $0.52$ ($\theta=1$). Each point rests on
  $R=20{,}000$ instances.}
  \label{fig:agents}
\end{figure}

Before evaluating the stored cells, we had formulated the hypothesis that more agents reduce the
EFX rate. We test it on all $360$ courses over $n=2,\dots,5$ at fixed $r$, $q$, agreement and
dispersion. In $284$ courses, the EFX rate falls strictly at every step. In $69$ courses, it is
constant because it stays at a bound: at zero under full agreement with $q>0$
(Proposition~\ref{prop:identical}), or at one for $q=0$ and Gini $0.06$ ($\theta=0.1$)
(Proposition~\ref{prop:theta-limit}). Five of the remaining seven courses are boundary cases that
drop to zero after $n=2$ or stay flat just below one. The other two show practically no effect
of the number of agents: for $r=3$, $q=1$, independent rankings and the two smallest levels of
dispersion, the EFX rate stays around $0.89$ to $0.91$ for every $n$. The decline also steepens
with agreement. With $q=0$, $r=3$ and Gini $0.52$, the EFX rate falls from $0.885$ to $0.690$
between two and five agents under independent rankings, but from $0.506$ to $0.047$ at Spearman
$0.73$ ($\rho=0.75$) (measured).

A natural reading is that the number of pairs who might envy each other grows quickly, from one
pair with two agents to ten pairs with five, while each pair alone does not become riskier. We
state this as a conjecture.

\subsection{Severity of Failures}\label{sec:agents-severity}

Figure~\ref{fig:agents-survival} shows the survival curves for two to five agents. A curve for
more agents never lies noticeably above a curve for fewer agents (by at most $0.0003$). At
$\alpha=0.2$ the curves are still close together, between $0.989$ and $0.972$, and from there on
they fan out. At $\alpha=0.6$, the share of instances that reach at least this level is
$0.870$, $0.756$, $0.651$ and $0.570$ for $n=2$ to $5$. Severe failures become somewhat more
frequent: the share of instances with $\astar<0.2$ rises from $0.011$ to $0.028$.

\begin{figure}[ht]
  \centering
  \includegraphics{fig45_agents_survival.pdf}
  \caption{Share of instances with $\astar\geq\alpha$ under Round Robin, one curve per number of
  agents, with $m=3n$ goods ($r=3$, $q=0$), Spearman $0.48$ ($\rho=0.5$) and Gini $0.52$
  ($\theta=1$). Each curve rests on $R=20{,}000$ instances. The point on each curve marks the
  $5\%$ quantile of $\astar$ among the instances that are not EFX.}
  \label{fig:agents-survival}
\end{figure}

Among the instances that fail, the picture is split. The conditional $5\%$ quantile hardly moves:
it is $0.23$ for every number of agents. The middle of the failures, however, moves down: the
median of $\astar$ among the instances that are not EFX falls from $0.66$ with two agents to
$0.64$, $0.60$ and $0.58$. Whether this happens depends on agreement. Under independent rankings
at the same dispersion, the median of the failures does not fall but stays between $0.72$ and
$0.74$, and the conditional $5\%$ quantile even rises slightly, from $0.28$ to $0.33$. At Spearman
$0.73$ ($\rho=0.75$), the median falls from $0.64$ to $0.50$. More agents therefore make failures
more frequent everywhere, and in a world with agreement they also make them deeper (measured).
The conditional $5\%$ quantile alone would have hidden the second part.

\paragraph{Relation to literature.}
We found no study of Round Robin with a growing number of agents at a fixed number of goods per
agent. The asymptotic results of \textcite{manurangsi2021} let the number of agents grow, but they
also let the number of goods per agent grow and assume independent values, so they do not speak
to our comparison. The existence of an EFX allocation is proven for two agents
\parencite{plaut2020} and three agents with additive valuations \parencite{chaudhury2024}; for
these cells, every failure is a failure of Round Robin. For four and five agents with $m=12$ and
$15$ goods, existence is an open problem \parencite[Open Problem~2]{amanatidis2023survey}, and we
assume it (Section~\ref{sec:existence}).

\paragraph{Recommendation.}
The more people take part, the less Round Robin can be relied on for EFX, even when everyone picks
equally often, and in a world where people agree, failures also become deeper. So far we have
judged Round Robin on its own; the next chapter compares it with what people would do without any
procedure.
